\documentclass[../../../include/open-logic-section]{subfiles}

\begin{document}

\olfileid{sth}{replacement}{extrinsic}
\olsection[Panglimbangan Ekstrinsik]{Panglimbangan Ekstrinsik babagan Panggantian}

Kita miwiti kanthi nimbang upaya
\emph{ekstrinsik} kanggo menehi dhasar marang
Panggantian. Boolos ngusulake mangkene:
\begin{quote}
  [\ldots] alasan nampa aksioma-aksioma panggantian
  iku cukup prasaja: akibat kang dikarepake akeh,
  dene akibat kang ora dikarepake (kaya katone)
  ora ana. Saliyane teorema babagan konsepsi
  iteratif, akibat-akibate kalebu teori wilangan
  tanpa wates kang nyukupi sanadyan durung ideal,
  lan asil kang banget dikarepake, kang menehi
  dhasar kanggo definisi induktif ing relasi kang
  saben bagean ora kosonge nduweni unsur minimal.
  \citep[229]{Boolos1971}
\end{quote}		
Intine pamanggih Boolos yaiku Panggantian kudu
diwenehi dhasar saka asil-asile. Asil kang
disebutake iku kang wis dirembug ing sawetara
bab sadurunge. Panggantian ngidini kita
mbuktekake yen ordinal von Neumann dadi wakil
kang apik banget kanggo gagasan tipe urutan
becik (iki kang diarani ``teori wilangan tanpa
wates kang nyukupi sanadyan durung ideal'').
Panggantian uga ngidini kita netepake
$V_\alpha$, netepake gagasan rank, lan
mbuktekake Induksi-$\in$ (iki cocog karo
``teorema babagan konsepsi iteratif'').
Pungkasan, Panggantian ngidini kita
mbuktekake Teorema Rekursi Transfinit (iki
``definisi induktif ing relasi kang saben bagean
ora kosonge nduweni unsur minimal'').

Iki pancen akibat kang dikarepake. Nanging,
apa akibat-akibat iki cukup kanggo
\emph{menehi dhasar} marang Panggantian?
\emph{Ora}. Paling ora, ora kanthi langsung.

Ana sawijining masalah prasaja. Sanadyan kita
wis nyebutake sawetara akibat Panggantian kang
dikarepake, akeh kang bisa dipikolehi kanthi
cara liya. Kasunyatan iki durung cukup kawentar,
mula perlu diterangake sedhela.

Ana teori himpunan kang prasaja, yaiku Teori
Tataran, utawa cekake $\LT$.\footnote{Versi-versi
awal $\LT$ diusulake dening \citet{Montague1965}
lan \citet{Scott1974}; banjur disederhanakake
lan dirembug sak buku dening \citet{Potter2004};
lan bubar \citet{ButtonLT1} luwih
nyederhanakake $\LT$.} Aksioma-aksioma $\LT$
mung Ekstensionalitas, Pamisahan, lan pratelan
yen saben himpunan iku himpunan bagean saka
sawijining \emph{tataran}; ``tataran'' ing kene
ditetepake kanthi premati supaya tataran-tataran
iku tumindake kaya $V_\alpha$ kang wis kita
kenal. Dadi $\ZF$ mbuktekake $\LT$; nanging
$\LT$ \emph{adoh} luwih ringkih tinimbang
$\ZF$. Malah, $\LT$ ora ngasilake Pasangan,
Himpunan Pangkat, Tanpa Wates, utawa Panggantian.
Tetepna $\Zr$ minangka asil nambahake Tanpa Wates
lan Himpunan Pangkat marang $\LT$; iki uga
ngasilake Pasangan, mula $\Zr$ paling ora padha
kuwate karo $\Z$. Nanging satemene $\Zr$
luwih kuwat tinimbang $\Z$, amarga nambahake
pratelan yen saben himpunan nduweni rank (mula
aku ngusulake jeneng $\Zr$). Pancen, $\Zr$
ngasilake teori ordinal kang nyukupi; asil-asil
kang mbagi hierarki dadi tataran-tataran kang
terurut becik; bukti Induksi-$\in$; lan
\emph{versi} Rekursi Transfinit.

Dadi, sanadyan Boolos durung ngerti bab iki,
kabeh akibat kang dikarepake lan disebutake
dening dheweke bisa dipikolehi \emph{tanpa}
Panggantian; cukup nganggo $\Zr$ tinimbang $\Z$.

(Yen mangkono, apa sebabe kita milih dalan
lumrah, yaiku mulang $\ZF$, dudu $\LT$ lan
$\Zr$? Ana rong alasan. Kapisan, amarga
sejarah, miwiti saka $\LT$ ora lumrah; kita
kepengin nyawisake pamaca supaya bisa maca
rembugan teori himpunan kang luwih baku.
Kapindho, yen wis siyap ngerteni $\LT$ lan
$\Zr$, pamaca bisa maca \citealt{Potter2004}
lan \citealt{ButtonLT1}.)

Mesthi wae, amarga $\Zr$ pancen luwih ringkih
tinimbang $\ZF$, ana asil kang bisa dibuktekake
ing $\ZF$ nanging durung diputusake dening
$\Zr$. Mula bisa wae wong nyoba menehi dhasar
ekstrinsik marang Panggantian kanthi nuduhake
salah siji asil mau. Nanging, yen wis ngerti
carane nggunakake $\Zr$, angel nemokake akeh
conto pratelan kang (a) diputusake dening
Panggantian nanging ora kanthi cara liya,
lan (b) miturut intuisi bener. (Kanggo katrangan
luwih akeh, delengen \citealt[\S13.2]{Potter2004}.)

Intine mangkene. Kanggo menehi dhasar
ekstrinsik kang ngyakinake marang Panggantian,
kudu ditemokake asil kang \emph{ora bisa}
dipikolehi tanpa Panggantian. Iki ora gampang.

Saiki ayo nimbang masalah liyane kang muncul
ing saben upaya menehi dhasar mung
sacara ekstrinsik marang Panggantian. (Masalah
iki mbokmenawa luwih dhasar tinimbang kang
kapisan.) Boolos ora mung nuduhake yen
Panggantian nduweni akeh akibat kang dikarepake.
Dheweke uga kandha yen Panggantian ``(kaya
katone) ora nduweni akibat kang ora dikarepake''.
Nanging, watesan ``kaya katone'' ing kurung iku
pancen wigati banget.

Elinga carane kita tekan kene: Komprehensi Naif
nemoni inkonsistensi, banjur kita nanggapi
kanthi nampa konsepsi himpunan kang kumulatif
lan iteratif. Konsepsi iki ngemot crita kang,
muga-muga, njamin konsistensine. Nanging yen
Panggantian ora bisa diwenehi dhasar saka crita
iku, kita (nganti saiki) ora duwe alasan
kanggo percaya yen $\ZF$ konsisten. Utawa luwih
persis: kita ora duwe alasan kanggo percaya
yen $\ZF$ konsisten kejaba kasunyatan (kang
mbokmenawa mung kebeneran) yen durung ana
wong nemokake kontradiksi \emph{nganti saiki}.
Ing pangertèn iku, pamanggih Boolos katon
padha karo iki: ``(kaya katone) $\ZF$
konsisten''. Kita kudune nuntut jaminan
konsistensi kang luwih kuwat tinimbang iki.

Masalah iki nyangkut saben upaya kang
\emph{mung} nganggo dhasar ekstrinsik kanggo
Panggantian, yaiku dhasar kang mung ngacu
marang akibat-akibat $\ZF$ kang wis dingerteni.
Mula kita bakal nggoleki dhasar
\emph{intrinsik} kanggo Panggantian: dhasar
kang nuduhake yen crita kita babagan himpunan
sejatine ``wis'' nuntun kita nampa Panggantian.

\end{document}
